%==============================================================================
% Section IV: Experiments and Results
%
% Included from main.tex via %==============================================================================
% Section IV: Experiments and Results
%
% Included from main.tex via %==============================================================================
% Section IV: Experiments and Results
%
% Included from main.tex via %==============================================================================
% Section IV: Experiments and Results
%
% Included from main.tex via \input{sections/results}
% Compile main.tex, not this file.
%==============================================================================
\section{Experiments and Results}
\label{sec:results}

%------------------------------------------------------------------------------
\subsection{Experimental Design}
\label{sec:design}

Four experiments were run, summarised in Table~\ref{tab:experiments}. The first
two correspond to the evaluation protocols of Section~\ref{sec:protocol} and
form the main results. The remaining two are ablations that remove one component
of the method at a time while holding everything else fixed. All four use the
same window length, stride, optimiser settings and epoch budget, so differences
between them are attributable to the stated change alone.

\begin{table}[!t]
\caption{Experiments. All use a window of 100 samples with stride 20, a maximum
of 100 epochs with early stopping, and a \SI{10}{\second} calibration window
unless stated otherwise.}
\label{tab:experiments}
\centering
\begin{tabular}{llp{3.0cm}}
\toprule
ID & Configuration & Question addressed \\
\midrule
E1 & Cross-day, four folds & Performance in the deployment setting \\
E2 & Within-day, three folds per session & Behaviour when between-session variation is removed \\
E3 & E1 without gain calibration & Contribution of calibration \\
E4 & E1 with a static base model & Contribution of base model dynamics \\
\bottomrule
\end{tabular}
\end{table}

All experiments were run on a single NVIDIA V100 GPU. Training the seven
architectures took 122 minutes for the four cross-day folds, 88 minutes for the
fifteen within-day folds, and 128 minutes for the dynamics ablation.

%------------------------------------------------------------------------------
\subsection{Cross-Day Performance}
\label{sec:crossday}

Table~\ref{tab:crossday} reports cross-day results. Each column is a session
held out entirely from training and validation, and each entry is the mean RMSE
over the recordings of that session. The row labelled EDM base model gives the
performance of the calibrated base model on exactly the same samples, so every
compensated result can be read against its own baseline.

\begin{table*}[!t]
\caption{Cross-day (leave-one-day-out) RMSE per held-out session. Lower is
better. The best compensated result in each column is shown in bold. Every
model improves on the base model in every fold.}
\label{tab:crossday}
\centering
\begin{tabular}{lrrrrrr}
\toprule
Model & 1st\_Comp & 2nd\_Comp & 3rd\_Base & 4th\_Valid & Overall RMSE & Overall MAE \\
\midrule
EDM base model          & 0.4884 & 1.0793 & 0.1890 & 0.3659 & 0.5581 & 0.4461 \\
\midrule
Transformer             & \textbf{0.3923} & 0.9310 & \textbf{0.1342} & \textbf{0.2450} & \textbf{0.4436} & 0.3252 \\
Conv-BiLSTM             & 0.4435 & 0.9138 & 0.1436 & 0.2480 & 0.4495 & 0.3310 \\
PINN                    & 0.4384 & 0.9063 & 0.1458 & 0.2557 & 0.4501 & 0.3363 \\
LSTM                    & 0.4536 & \textbf{0.8915} & 0.1488 & 0.2638 & 0.4520 & 0.3314 \\
Conv-BiLSTM-Transformer & 0.4482 & 0.9258 & 0.1396 & 0.2521 & 0.4547 & 0.3371 \\
Selective SSM (Mamba)   & 0.4716 & 0.9010 & 0.1494 & 0.2750 & 0.4622 & 0.3448 \\
Synergy-SSM             & 0.4637 & 0.8972 & 0.1574 & 0.2830 & 0.4644 & 0.3441 \\
\bottomrule
\end{tabular}
\end{table*}

The compensator improves on the base model in every fold and for every
architecture. Taken over all held-out recordings, the best model reduces RMSE
from 0.5581 to 0.4436, a relative reduction of 20.5 percent, and even the
weakest of the seven reaches 0.4644, a reduction of 16.8 percent. Of the 224
model and recording combinations, 175 (78 percent) show lower error than the
base model on the same samples.

The size of the improvement varies with the held-out session. It is smallest for
2nd\_Comp (17.4 percent), which has the largest force range and the highest gain
of the four, and largest for 4th\_Valid (33.0 percent) and 3rd\_Base
(29.0 percent), whose gains sit near the middle of the range spanned by the
training sessions. The improvement on 1st\_Comp is 19.7 percent. This ordering is
consistent with the interpretation developed in Section~\ref{sec:extrapolation}:
the method performs best when the held-out session is well inside the range of
conditions represented in training.

A one-sided reading of these numbers would be unwarranted without testing them,
so we treat each held-out recording as a paired observation and compare each
model against the base model with a Wilcoxon signed-rank test over the 32 test
recordings. All seven architectures improve on the base model at
$p < 4 \times 10^{-4}$, with the strongest evidence for the PINN
($p = 5.9 \times 10^{-7}$, improving on 28 of 32 recordings) and the Transformer
and Conv-BiLSTM ($p = 1.3 \times 10^{-5}$, 26 of 32 each). The benefit of
compensation is therefore consistent across recordings rather than driven by a
few favourable cases.

%------------------------------------------------------------------------------
\subsection{Within-Day Performance}
\label{sec:withinday}

Table~\ref{tab:withinday} reports the within-day results, in which training and
testing use different recordings from the same session. Across all 36 held-out
recordings the picture is markedly weaker than in the cross-day setting. Only
two of the seven architectures achieve a lower aggregate RMSE than the base
model, and the best of them improves on it by 2.9 percent, compared with
20.5 percent across sessions.

\begin{table}[!t]
\caption{Within-day RMSE, averaged over three rotating folds per session. The
base model is placed in rank order so that models falling below it are visible
directly.}
\label{tab:withinday}
\centering
\begin{tabular}{lrr}
\toprule
Model & Overall RMSE & Overall MAE \\
\midrule
PINN                    & \textbf{0.4343} & 0.3354 \\
Synergy-SSM             & 0.4467 & 0.3474 \\
\textit{EDM base model} & \textit{0.4473} & \textit{0.3562} \\
Conv-BiLSTM             & 0.4481 & 0.3513 \\
Selective SSM (Mamba)   & 0.4498 & 0.3593 \\
Transformer             & 0.4610 & 0.3645 \\
LSTM                    & 0.4624 & 0.3719 \\
Conv-BiLSTM-Transformer & 0.4633 & 0.3673 \\
\bottomrule
\end{tabular}
\end{table}

This aggregate requires careful interpretation, because it is not representative
of the individual sessions. Averaging RMSE over recordings whose force ranges
differ by more than an order of magnitude gives disproportionate weight to the
session with the largest absolute errors. Here that session is 2nd\_Comp, which
contributes nine of the 36 held-out recordings and has errors near 1.1, roughly
twenty times those of 5th\_Comp. The PINN is the only architecture that improves
on the base model on 2nd\_Comp, which is the main reason it leads the aggregate.

Examined session by session, the result is more favourable. In every one of the
five sessions at least one architecture achieves a lower RMSE than the base
model: by 7.9 percent on 1st\_Comp, 3.6 percent on 2nd\_Comp, 13.2 percent on
3rd\_Base, 13.9 percent on 4th\_Valid and 11.5 percent on 5th\_Comp. What does
not hold is consistency. The architecture that performs best differs between
sessions, and no architecture improves on the base model in all five. Over all
model and recording combinations, 168 of 252 (67 percent) improve on the base
model, compared with 78 percent in the cross-day setting.

The contrast with Section~\ref{sec:crossday} is the main finding of this
experiment, and it runs opposite to the usual expectation that within-session
evaluation is the easier case. The explanation is the amount of training data
available. A cross-day fold trains on 24 to 28 recordings, whereas a within-day
fold trains on between three and nine recordings from a single session. Removing
between-session variation also removes most of the training set, and for this
task the loss of data outweighs the gain in homogeneity. A secondary factor is
that per-session calibration (Section~\ref{sec:calibration}) already absorbs much
of the between-session gain variation before the network sees the data, so the
cross-day setting is less disadvantaged than it would otherwise be.

%------------------------------------------------------------------------------
\subsection{Ablation: Per-Session Gain Calibration}
\label{sec:ablation-calib}

% PENDING. E3 did not complete: two concurrently submitted jobs wrote to the
% same output directory, and the first to finish relocated the directory while
% E3 was still running. The run must be repeated in isolation before this
% subsection can be written.
%
% Expected content once available:
%   - cross-day RMSE for base and all seven models with calibration disabled
%   - comparison against Table IV to isolate the contribution of calibration
%   - this supports the paper's central methodological claim, so it should be
%     reported with the same per-fold detail as Section IV-B

%------------------------------------------------------------------------------
\subsection{Ablation: Base Model Dynamics}
\label{sec:ablation-lpf}

Section~\ref{sec:basemodel} argued that the first-order stage is required
because the force to sEMG relationship is not single-valued. To test this, the
cross-day experiment was repeated with the first-order stage removed, leaving a
purely static saturating nonlinearity. All other settings were unchanged.

\begin{table}[!t]
\caption{Effect of removing the first-order stage from the base model. Values
are cross-day RMSE over all held-out recordings.}
\label{tab:lpf}
\centering
\begin{tabular}{lrr}
\toprule
 & With dynamics & Static only \\
\midrule
EDM base model          & 0.5581 & 0.5683 \\
Transformer             & \textbf{0.4436} & 0.4665 \\
Conv-BiLSTM             & 0.4495 & 0.4687 \\
Conv-BiLSTM-Transformer & 0.4547 & 0.4697 \\
PINN                    & 0.4501 & 0.4845 \\
LSTM                    & 0.4520 & 0.4913 \\
Synergy-SSM             & 0.4644 & 0.4917 \\
Selective SSM (Mamba)   & 0.4622 & 0.4929 \\
\bottomrule
\end{tabular}
\end{table}

Removing the dynamics degrades both stages of the method, as shown in
Table~\ref{tab:lpf}. The base model itself worsens from 0.5581 to 0.5683, and
the best compensated result worsens from 0.4436 to 0.4665, a relative increase
in error of 4.9 percent. The degradation is consistent rather than selective:
all seven architectures perform worse with the static base model, and the number
of model and recording combinations improving on the base model falls from 175
to 164 out of 224.

Two observations follow. First, the compensator does not recover what the static
base model loses, even though it has access to a 100-sample history and could in
principle represent the missing dynamics itself. Placing the dynamics in the
physics-based component is therefore more effective than leaving them to be
learned. Second, the ranking of architectures is largely preserved between the
two conditions, which suggests that the choice of architecture and the choice of
base model structure act independently.

%------------------------------------------------------------------------------
\subsection{Transfer Boundary}
\label{sec:extrapolation}

Section~\ref{sec:protocol} explained that one session is reserved for training
because its gain lies well below that of all others. Treating it as a held-out
test session instead is informative, because it shows what happens when the
method is asked to extrapolate rather than interpolate. We therefore repeated
the cross-day experiment with that session held out, keeping all other settings
identical.

The outcome reverses the main result. The calibrated base model achieves an RMSE
of 0.0663 on this session, and every one of the seven architectures performs
worse, ranging from 0.0772 for the Transformer to 0.0838 for the Synergy-SSM.
The best model is thus 16.4 percent worse than applying no compensation at all,
and only 14 of 70 model and recording combinations improve on the base model,
compared with 78 percent in Table~\ref{tab:crossday}.

The distinguishing feature of this session is that its gain of 0.44 lies a factor
of 1.9 below the lowest gain among the remaining sessions, so no training session
brackets it from below. The compensator has learned a correction appropriate to
the range of conditions it observed and applies it outside that range, where it
is systematically wrong. This delimits the applicability of the method: it
improves accuracy when the deployment session falls inside the range of
conditions spanned by training, and degrades accuracy when it does not. Because
the gain can be estimated from the calibration segment before any prediction is
made, this condition is checkable at deployment time, and falling back to the
uncompensated base model is a safe response when it is not met.

%------------------------------------------------------------------------------
\subsection{Comparison Between Architectures}
\label{sec:architecture}

The seven architectures span convolutional, recurrent, attention-based,
state-space and physics-constrained designs, and differ by a factor of five in
parameter count, from 30{,}770 to 155{,}297. Despite this, their cross-day
results occupy a narrow band: mean RMSE ranges from 0.4436 to 0.4644, a relative
spread of 4.7 percent, against a 20.5 percent improvement over the base model.
The difference between the best and worst architecture is therefore about a
quarter of the difference between using compensation and not using it.

We tested this directly. A Friedman test across the 32 held-out recordings
rejects the hypothesis that the seven architectures have identical error
distributions ($\chi^2 = 23.88$, $p = 5 \times 10^{-4}$), so the models are not
strictly interchangeable. However, when all 21 pairwise comparisons are made with
Wilcoxon signed-rank tests and corrected for multiple comparisons using the Holm
procedure, no pair differs significantly at the 5 percent level. The smallest
adjusted $p$ value is 0.221, for LSTM against the selective state-space model.

The appropriate conclusion is therefore narrower than architectural equivalence.
There is detectable heterogeneity among the seven, but the data do not identify
which architecture is preferable: no individual comparison survives correction,
and the spread is small relative to the effect of compensation itself. For this
task and dataset, the choice of sequence architecture is not the factor that
determines accuracy. Parameter count in particular is uninformative, since the
largest model ranks second and the smallest ranks sixth.

This has a practical consequence. Where two architectures cannot be separated on
accuracy, secondary criteria become decisive. The selective state-space model
attains comparable error with the fewest parameters but the longest training
time, owing to its sequential scan. The LSTM is the only strictly causal model in
the comparison and is therefore the natural candidate for a streaming
implementation. The Transformer achieves the lowest mean error and trains
quickly, but it processes the whole window bidirectionally and is less suited to
real-time use.

% Compile main.tex, not this file.
%==============================================================================
\section{Experiments and Results}
\label{sec:results}

%------------------------------------------------------------------------------
\subsection{Experimental Design}
\label{sec:design}

Four experiments were run, summarised in Table~\ref{tab:experiments}. The first
two correspond to the evaluation protocols of Section~\ref{sec:protocol} and
form the main results. The remaining two are ablations that remove one component
of the method at a time while holding everything else fixed. All four use the
same window length, stride, optimiser settings and epoch budget, so differences
between them are attributable to the stated change alone.

\begin{table}[!t]
\caption{Experiments. All use a window of 100 samples with stride 20, a maximum
of 100 epochs with early stopping, and a \SI{10}{\second} calibration window
unless stated otherwise.}
\label{tab:experiments}
\centering
\begin{tabular}{llp{3.0cm}}
\toprule
ID & Configuration & Question addressed \\
\midrule
E1 & Cross-day, four folds & Performance in the deployment setting \\
E2 & Within-day, three folds per session & Behaviour when between-session variation is removed \\
E3 & E1 without gain calibration & Contribution of calibration \\
E4 & E1 with a static base model & Contribution of base model dynamics \\
\bottomrule
\end{tabular}
\end{table}

All experiments were run on a single NVIDIA V100 GPU. Training the seven
architectures took 122 minutes for the four cross-day folds, 88 minutes for the
fifteen within-day folds, and 128 minutes for the dynamics ablation.

%------------------------------------------------------------------------------
\subsection{Cross-Day Performance}
\label{sec:crossday}

Table~\ref{tab:crossday} reports cross-day results. Each column is a session
held out entirely from training and validation, and each entry is the mean RMSE
over the recordings of that session. The row labelled EDM base model gives the
performance of the calibrated base model on exactly the same samples, so every
compensated result can be read against its own baseline.

\begin{table*}[!t]
\caption{Cross-day (leave-one-day-out) RMSE per held-out session. Lower is
better. The best compensated result in each column is shown in bold. Every
model improves on the base model in every fold.}
\label{tab:crossday}
\centering
\begin{tabular}{lrrrrrr}
\toprule
Model & 1st\_Comp & 2nd\_Comp & 3rd\_Base & 4th\_Valid & Overall RMSE & Overall MAE \\
\midrule
EDM base model          & 0.4884 & 1.0793 & 0.1890 & 0.3659 & 0.5581 & 0.4461 \\
\midrule
Transformer             & \textbf{0.3923} & 0.9310 & \textbf{0.1342} & \textbf{0.2450} & \textbf{0.4436} & 0.3252 \\
Conv-BiLSTM             & 0.4435 & 0.9138 & 0.1436 & 0.2480 & 0.4495 & 0.3310 \\
PINN                    & 0.4384 & 0.9063 & 0.1458 & 0.2557 & 0.4501 & 0.3363 \\
LSTM                    & 0.4536 & \textbf{0.8915} & 0.1488 & 0.2638 & 0.4520 & 0.3314 \\
Conv-BiLSTM-Transformer & 0.4482 & 0.9258 & 0.1396 & 0.2521 & 0.4547 & 0.3371 \\
Selective SSM (Mamba)   & 0.4716 & 0.9010 & 0.1494 & 0.2750 & 0.4622 & 0.3448 \\
Synergy-SSM             & 0.4637 & 0.8972 & 0.1574 & 0.2830 & 0.4644 & 0.3441 \\
\bottomrule
\end{tabular}
\end{table*}

The compensator improves on the base model in every fold and for every
architecture. Taken over all held-out recordings, the best model reduces RMSE
from 0.5581 to 0.4436, a relative reduction of 20.5 percent, and even the
weakest of the seven reaches 0.4644, a reduction of 16.8 percent. Of the 224
model and recording combinations, 175 (78 percent) show lower error than the
base model on the same samples.

The size of the improvement varies with the held-out session. It is smallest for
2nd\_Comp (17.4 percent), which has the largest force range and the highest gain
of the four, and largest for 4th\_Valid (33.0 percent) and 3rd\_Base
(29.0 percent), whose gains sit near the middle of the range spanned by the
training sessions. The improvement on 1st\_Comp is 19.7 percent. This ordering is
consistent with the interpretation developed in Section~\ref{sec:extrapolation}:
the method performs best when the held-out session is well inside the range of
conditions represented in training.

A one-sided reading of these numbers would be unwarranted without testing them,
so we treat each held-out recording as a paired observation and compare each
model against the base model with a Wilcoxon signed-rank test over the 32 test
recordings. All seven architectures improve on the base model at
$p < 4 \times 10^{-4}$, with the strongest evidence for the PINN
($p = 5.9 \times 10^{-7}$, improving on 28 of 32 recordings) and the Transformer
and Conv-BiLSTM ($p = 1.3 \times 10^{-5}$, 26 of 32 each). The benefit of
compensation is therefore consistent across recordings rather than driven by a
few favourable cases.

%------------------------------------------------------------------------------
\subsection{Within-Day Performance}
\label{sec:withinday}

Table~\ref{tab:withinday} reports the within-day results, in which training and
testing use different recordings from the same session. Across all 36 held-out
recordings the picture is markedly weaker than in the cross-day setting. Only
two of the seven architectures achieve a lower aggregate RMSE than the base
model, and the best of them improves on it by 2.9 percent, compared with
20.5 percent across sessions.

\begin{table}[!t]
\caption{Within-day RMSE, averaged over three rotating folds per session. The
base model is placed in rank order so that models falling below it are visible
directly.}
\label{tab:withinday}
\centering
\begin{tabular}{lrr}
\toprule
Model & Overall RMSE & Overall MAE \\
\midrule
PINN                    & \textbf{0.4343} & 0.3354 \\
Synergy-SSM             & 0.4467 & 0.3474 \\
\textit{EDM base model} & \textit{0.4473} & \textit{0.3562} \\
Conv-BiLSTM             & 0.4481 & 0.3513 \\
Selective SSM (Mamba)   & 0.4498 & 0.3593 \\
Transformer             & 0.4610 & 0.3645 \\
LSTM                    & 0.4624 & 0.3719 \\
Conv-BiLSTM-Transformer & 0.4633 & 0.3673 \\
\bottomrule
\end{tabular}
\end{table}

This aggregate requires careful interpretation, because it is not representative
of the individual sessions. Averaging RMSE over recordings whose force ranges
differ by more than an order of magnitude gives disproportionate weight to the
session with the largest absolute errors. Here that session is 2nd\_Comp, which
contributes nine of the 36 held-out recordings and has errors near 1.1, roughly
twenty times those of 5th\_Comp. The PINN is the only architecture that improves
on the base model on 2nd\_Comp, which is the main reason it leads the aggregate.

Examined session by session, the result is more favourable. In every one of the
five sessions at least one architecture achieves a lower RMSE than the base
model: by 7.9 percent on 1st\_Comp, 3.6 percent on 2nd\_Comp, 13.2 percent on
3rd\_Base, 13.9 percent on 4th\_Valid and 11.5 percent on 5th\_Comp. What does
not hold is consistency. The architecture that performs best differs between
sessions, and no architecture improves on the base model in all five. Over all
model and recording combinations, 168 of 252 (67 percent) improve on the base
model, compared with 78 percent in the cross-day setting.

The contrast with Section~\ref{sec:crossday} is the main finding of this
experiment, and it runs opposite to the usual expectation that within-session
evaluation is the easier case. The explanation is the amount of training data
available. A cross-day fold trains on 24 to 28 recordings, whereas a within-day
fold trains on between three and nine recordings from a single session. Removing
between-session variation also removes most of the training set, and for this
task the loss of data outweighs the gain in homogeneity. A secondary factor is
that per-session calibration (Section~\ref{sec:calibration}) already absorbs much
of the between-session gain variation before the network sees the data, so the
cross-day setting is less disadvantaged than it would otherwise be.

%------------------------------------------------------------------------------
\subsection{Ablation: Per-Session Gain Calibration}
\label{sec:ablation-calib}

% PENDING. E3 did not complete: two concurrently submitted jobs wrote to the
% same output directory, and the first to finish relocated the directory while
% E3 was still running. The run must be repeated in isolation before this
% subsection can be written.
%
% Expected content once available:
%   - cross-day RMSE for base and all seven models with calibration disabled
%   - comparison against Table IV to isolate the contribution of calibration
%   - this supports the paper's central methodological claim, so it should be
%     reported with the same per-fold detail as Section IV-B

%------------------------------------------------------------------------------
\subsection{Ablation: Base Model Dynamics}
\label{sec:ablation-lpf}

Section~\ref{sec:basemodel} argued that the first-order stage is required
because the force to sEMG relationship is not single-valued. To test this, the
cross-day experiment was repeated with the first-order stage removed, leaving a
purely static saturating nonlinearity. All other settings were unchanged.

\begin{table}[!t]
\caption{Effect of removing the first-order stage from the base model. Values
are cross-day RMSE over all held-out recordings.}
\label{tab:lpf}
\centering
\begin{tabular}{lrr}
\toprule
 & With dynamics & Static only \\
\midrule
EDM base model          & 0.5581 & 0.5683 \\
Transformer             & \textbf{0.4436} & 0.4665 \\
Conv-BiLSTM             & 0.4495 & 0.4687 \\
Conv-BiLSTM-Transformer & 0.4547 & 0.4697 \\
PINN                    & 0.4501 & 0.4845 \\
LSTM                    & 0.4520 & 0.4913 \\
Synergy-SSM             & 0.4644 & 0.4917 \\
Selective SSM (Mamba)   & 0.4622 & 0.4929 \\
\bottomrule
\end{tabular}
\end{table}

Removing the dynamics degrades both stages of the method, as shown in
Table~\ref{tab:lpf}. The base model itself worsens from 0.5581 to 0.5683, and
the best compensated result worsens from 0.4436 to 0.4665, a relative increase
in error of 4.9 percent. The degradation is consistent rather than selective:
all seven architectures perform worse with the static base model, and the number
of model and recording combinations improving on the base model falls from 175
to 164 out of 224.

Two observations follow. First, the compensator does not recover what the static
base model loses, even though it has access to a 100-sample history and could in
principle represent the missing dynamics itself. Placing the dynamics in the
physics-based component is therefore more effective than leaving them to be
learned. Second, the ranking of architectures is largely preserved between the
two conditions, which suggests that the choice of architecture and the choice of
base model structure act independently.

%------------------------------------------------------------------------------
\subsection{Transfer Boundary}
\label{sec:extrapolation}

Section~\ref{sec:protocol} explained that one session is reserved for training
because its gain lies well below that of all others. Treating it as a held-out
test session instead is informative, because it shows what happens when the
method is asked to extrapolate rather than interpolate. We therefore repeated
the cross-day experiment with that session held out, keeping all other settings
identical.

The outcome reverses the main result. The calibrated base model achieves an RMSE
of 0.0663 on this session, and every one of the seven architectures performs
worse, ranging from 0.0772 for the Transformer to 0.0838 for the Synergy-SSM.
The best model is thus 16.4 percent worse than applying no compensation at all,
and only 14 of 70 model and recording combinations improve on the base model,
compared with 78 percent in Table~\ref{tab:crossday}.

The distinguishing feature of this session is that its gain of 0.44 lies a factor
of 1.9 below the lowest gain among the remaining sessions, so no training session
brackets it from below. The compensator has learned a correction appropriate to
the range of conditions it observed and applies it outside that range, where it
is systematically wrong. This delimits the applicability of the method: it
improves accuracy when the deployment session falls inside the range of
conditions spanned by training, and degrades accuracy when it does not. Because
the gain can be estimated from the calibration segment before any prediction is
made, this condition is checkable at deployment time, and falling back to the
uncompensated base model is a safe response when it is not met.

%------------------------------------------------------------------------------
\subsection{Comparison Between Architectures}
\label{sec:architecture}

The seven architectures span convolutional, recurrent, attention-based,
state-space and physics-constrained designs, and differ by a factor of five in
parameter count, from 30{,}770 to 155{,}297. Despite this, their cross-day
results occupy a narrow band: mean RMSE ranges from 0.4436 to 0.4644, a relative
spread of 4.7 percent, against a 20.5 percent improvement over the base model.
The difference between the best and worst architecture is therefore about a
quarter of the difference between using compensation and not using it.

We tested this directly. A Friedman test across the 32 held-out recordings
rejects the hypothesis that the seven architectures have identical error
distributions ($\chi^2 = 23.88$, $p = 5 \times 10^{-4}$), so the models are not
strictly interchangeable. However, when all 21 pairwise comparisons are made with
Wilcoxon signed-rank tests and corrected for multiple comparisons using the Holm
procedure, no pair differs significantly at the 5 percent level. The smallest
adjusted $p$ value is 0.221, for LSTM against the selective state-space model.

The appropriate conclusion is therefore narrower than architectural equivalence.
There is detectable heterogeneity among the seven, but the data do not identify
which architecture is preferable: no individual comparison survives correction,
and the spread is small relative to the effect of compensation itself. For this
task and dataset, the choice of sequence architecture is not the factor that
determines accuracy. Parameter count in particular is uninformative, since the
largest model ranks second and the smallest ranks sixth.

This has a practical consequence. Where two architectures cannot be separated on
accuracy, secondary criteria become decisive. The selective state-space model
attains comparable error with the fewest parameters but the longest training
time, owing to its sequential scan. The LSTM is the only strictly causal model in
the comparison and is therefore the natural candidate for a streaming
implementation. The Transformer achieves the lowest mean error and trains
quickly, but it processes the whole window bidirectionally and is less suited to
real-time use.

% Compile main.tex, not this file.
%==============================================================================
\section{Experiments and Results}
\label{sec:results}

%------------------------------------------------------------------------------
\subsection{Experimental Design}
\label{sec:design}

Four experiments were run, summarised in Table~\ref{tab:experiments}. The first
two correspond to the evaluation protocols of Section~\ref{sec:protocol} and
form the main results. The remaining two are ablations that remove one component
of the method at a time while holding everything else fixed. All four use the
same window length, stride, optimiser settings and epoch budget, so differences
between them are attributable to the stated change alone.

\begin{table}[!t]
\caption{Experiments. All use a window of 100 samples with stride 20, a maximum
of 100 epochs with early stopping, and a \SI{10}{\second} calibration window
unless stated otherwise.}
\label{tab:experiments}
\centering
\begin{tabular}{llp{3.0cm}}
\toprule
ID & Configuration & Question addressed \\
\midrule
E1 & Cross-day, four folds & Performance in the deployment setting \\
E2 & Within-day, three folds per session & Behaviour when between-session variation is removed \\
E3 & E1 without gain calibration & Contribution of calibration \\
E4 & E1 with a static base model & Contribution of base model dynamics \\
\bottomrule
\end{tabular}
\end{table}

All experiments were run on a single NVIDIA V100 GPU. Training the seven
architectures took 122 minutes for the four cross-day folds, 88 minutes for the
fifteen within-day folds, and 128 minutes for the dynamics ablation.

%------------------------------------------------------------------------------
\subsection{Cross-Day Performance}
\label{sec:crossday}

Table~\ref{tab:crossday} reports cross-day results. Each column is a session
held out entirely from training and validation, and each entry is the mean RMSE
over the recordings of that session. The row labelled EDM base model gives the
performance of the calibrated base model on exactly the same samples, so every
compensated result can be read against its own baseline.

\begin{table*}[!t]
\caption{Cross-day (leave-one-day-out) RMSE per held-out session. Lower is
better. The best compensated result in each column is shown in bold. Every
model improves on the base model in every fold.}
\label{tab:crossday}
\centering
\begin{tabular}{lrrrrrr}
\toprule
Model & 1st\_Comp & 2nd\_Comp & 3rd\_Base & 4th\_Valid & Overall RMSE & Overall MAE \\
\midrule
EDM base model          & 0.4884 & 1.0793 & 0.1890 & 0.3659 & 0.5581 & 0.4461 \\
\midrule
Transformer             & \textbf{0.3923} & 0.9310 & \textbf{0.1342} & \textbf{0.2450} & \textbf{0.4436} & 0.3252 \\
Conv-BiLSTM             & 0.4435 & 0.9138 & 0.1436 & 0.2480 & 0.4495 & 0.3310 \\
PINN                    & 0.4384 & 0.9063 & 0.1458 & 0.2557 & 0.4501 & 0.3363 \\
LSTM                    & 0.4536 & \textbf{0.8915} & 0.1488 & 0.2638 & 0.4520 & 0.3314 \\
Conv-BiLSTM-Transformer & 0.4482 & 0.9258 & 0.1396 & 0.2521 & 0.4547 & 0.3371 \\
Selective SSM (Mamba)   & 0.4716 & 0.9010 & 0.1494 & 0.2750 & 0.4622 & 0.3448 \\
Synergy-SSM             & 0.4637 & 0.8972 & 0.1574 & 0.2830 & 0.4644 & 0.3441 \\
\bottomrule
\end{tabular}
\end{table*}

The compensator improves on the base model in every fold and for every
architecture. Taken over all held-out recordings, the best model reduces RMSE
from 0.5581 to 0.4436, a relative reduction of 20.5 percent, and even the
weakest of the seven reaches 0.4644, a reduction of 16.8 percent. Of the 224
model and recording combinations, 175 (78 percent) show lower error than the
base model on the same samples.

The size of the improvement varies with the held-out session. It is smallest for
2nd\_Comp (17.4 percent), which has the largest force range and the highest gain
of the four, and largest for 4th\_Valid (33.0 percent) and 3rd\_Base
(29.0 percent), whose gains sit near the middle of the range spanned by the
training sessions. The improvement on 1st\_Comp is 19.7 percent. This ordering is
consistent with the interpretation developed in Section~\ref{sec:extrapolation}:
the method performs best when the held-out session is well inside the range of
conditions represented in training.

A one-sided reading of these numbers would be unwarranted without testing them,
so we treat each held-out recording as a paired observation and compare each
model against the base model with a Wilcoxon signed-rank test over the 32 test
recordings. All seven architectures improve on the base model at
$p < 4 \times 10^{-4}$, with the strongest evidence for the PINN
($p = 5.9 \times 10^{-7}$, improving on 28 of 32 recordings) and the Transformer
and Conv-BiLSTM ($p = 1.3 \times 10^{-5}$, 26 of 32 each). The benefit of
compensation is therefore consistent across recordings rather than driven by a
few favourable cases.

%------------------------------------------------------------------------------
\subsection{Within-Day Performance}
\label{sec:withinday}

Table~\ref{tab:withinday} reports the within-day results, in which training and
testing use different recordings from the same session. Across all 36 held-out
recordings the picture is markedly weaker than in the cross-day setting. Only
two of the seven architectures achieve a lower aggregate RMSE than the base
model, and the best of them improves on it by 2.9 percent, compared with
20.5 percent across sessions.

\begin{table}[!t]
\caption{Within-day RMSE, averaged over three rotating folds per session. The
base model is placed in rank order so that models falling below it are visible
directly.}
\label{tab:withinday}
\centering
\begin{tabular}{lrr}
\toprule
Model & Overall RMSE & Overall MAE \\
\midrule
PINN                    & \textbf{0.4343} & 0.3354 \\
Synergy-SSM             & 0.4467 & 0.3474 \\
\textit{EDM base model} & \textit{0.4473} & \textit{0.3562} \\
Conv-BiLSTM             & 0.4481 & 0.3513 \\
Selective SSM (Mamba)   & 0.4498 & 0.3593 \\
Transformer             & 0.4610 & 0.3645 \\
LSTM                    & 0.4624 & 0.3719 \\
Conv-BiLSTM-Transformer & 0.4633 & 0.3673 \\
\bottomrule
\end{tabular}
\end{table}

This aggregate requires careful interpretation, because it is not representative
of the individual sessions. Averaging RMSE over recordings whose force ranges
differ by more than an order of magnitude gives disproportionate weight to the
session with the largest absolute errors. Here that session is 2nd\_Comp, which
contributes nine of the 36 held-out recordings and has errors near 1.1, roughly
twenty times those of 5th\_Comp. The PINN is the only architecture that improves
on the base model on 2nd\_Comp, which is the main reason it leads the aggregate.

Examined session by session, the result is more favourable. In every one of the
five sessions at least one architecture achieves a lower RMSE than the base
model: by 7.9 percent on 1st\_Comp, 3.6 percent on 2nd\_Comp, 13.2 percent on
3rd\_Base, 13.9 percent on 4th\_Valid and 11.5 percent on 5th\_Comp. What does
not hold is consistency. The architecture that performs best differs between
sessions, and no architecture improves on the base model in all five. Over all
model and recording combinations, 168 of 252 (67 percent) improve on the base
model, compared with 78 percent in the cross-day setting.

The contrast with Section~\ref{sec:crossday} is the main finding of this
experiment, and it runs opposite to the usual expectation that within-session
evaluation is the easier case. The explanation is the amount of training data
available. A cross-day fold trains on 24 to 28 recordings, whereas a within-day
fold trains on between three and nine recordings from a single session. Removing
between-session variation also removes most of the training set, and for this
task the loss of data outweighs the gain in homogeneity. A secondary factor is
that per-session calibration (Section~\ref{sec:calibration}) already absorbs much
of the between-session gain variation before the network sees the data, so the
cross-day setting is less disadvantaged than it would otherwise be.

%------------------------------------------------------------------------------
\subsection{Ablation: Per-Session Gain Calibration}
\label{sec:ablation-calib}

% PENDING. E3 did not complete: two concurrently submitted jobs wrote to the
% same output directory, and the first to finish relocated the directory while
% E3 was still running. The run must be repeated in isolation before this
% subsection can be written.
%
% Expected content once available:
%   - cross-day RMSE for base and all seven models with calibration disabled
%   - comparison against Table IV to isolate the contribution of calibration
%   - this supports the paper's central methodological claim, so it should be
%     reported with the same per-fold detail as Section IV-B

%------------------------------------------------------------------------------
\subsection{Ablation: Base Model Dynamics}
\label{sec:ablation-lpf}

Section~\ref{sec:basemodel} argued that the first-order stage is required
because the force to sEMG relationship is not single-valued. To test this, the
cross-day experiment was repeated with the first-order stage removed, leaving a
purely static saturating nonlinearity. All other settings were unchanged.

\begin{table}[!t]
\caption{Effect of removing the first-order stage from the base model. Values
are cross-day RMSE over all held-out recordings.}
\label{tab:lpf}
\centering
\begin{tabular}{lrr}
\toprule
 & With dynamics & Static only \\
\midrule
EDM base model          & 0.5581 & 0.5683 \\
Transformer             & \textbf{0.4436} & 0.4665 \\
Conv-BiLSTM             & 0.4495 & 0.4687 \\
Conv-BiLSTM-Transformer & 0.4547 & 0.4697 \\
PINN                    & 0.4501 & 0.4845 \\
LSTM                    & 0.4520 & 0.4913 \\
Synergy-SSM             & 0.4644 & 0.4917 \\
Selective SSM (Mamba)   & 0.4622 & 0.4929 \\
\bottomrule
\end{tabular}
\end{table}

Removing the dynamics degrades both stages of the method, as shown in
Table~\ref{tab:lpf}. The base model itself worsens from 0.5581 to 0.5683, and
the best compensated result worsens from 0.4436 to 0.4665, a relative increase
in error of 4.9 percent. The degradation is consistent rather than selective:
all seven architectures perform worse with the static base model, and the number
of model and recording combinations improving on the base model falls from 175
to 164 out of 224.

Two observations follow. First, the compensator does not recover what the static
base model loses, even though it has access to a 100-sample history and could in
principle represent the missing dynamics itself. Placing the dynamics in the
physics-based component is therefore more effective than leaving them to be
learned. Second, the ranking of architectures is largely preserved between the
two conditions, which suggests that the choice of architecture and the choice of
base model structure act independently.

%------------------------------------------------------------------------------
\subsection{Transfer Boundary}
\label{sec:extrapolation}

Section~\ref{sec:protocol} explained that one session is reserved for training
because its gain lies well below that of all others. Treating it as a held-out
test session instead is informative, because it shows what happens when the
method is asked to extrapolate rather than interpolate. We therefore repeated
the cross-day experiment with that session held out, keeping all other settings
identical.

The outcome reverses the main result. The calibrated base model achieves an RMSE
of 0.0663 on this session, and every one of the seven architectures performs
worse, ranging from 0.0772 for the Transformer to 0.0838 for the Synergy-SSM.
The best model is thus 16.4 percent worse than applying no compensation at all,
and only 14 of 70 model and recording combinations improve on the base model,
compared with 78 percent in Table~\ref{tab:crossday}.

The distinguishing feature of this session is that its gain of 0.44 lies a factor
of 1.9 below the lowest gain among the remaining sessions, so no training session
brackets it from below. The compensator has learned a correction appropriate to
the range of conditions it observed and applies it outside that range, where it
is systematically wrong. This delimits the applicability of the method: it
improves accuracy when the deployment session falls inside the range of
conditions spanned by training, and degrades accuracy when it does not. Because
the gain can be estimated from the calibration segment before any prediction is
made, this condition is checkable at deployment time, and falling back to the
uncompensated base model is a safe response when it is not met.

%------------------------------------------------------------------------------
\subsection{Comparison Between Architectures}
\label{sec:architecture}

The seven architectures span convolutional, recurrent, attention-based,
state-space and physics-constrained designs, and differ by a factor of five in
parameter count, from 30{,}770 to 155{,}297. Despite this, their cross-day
results occupy a narrow band: mean RMSE ranges from 0.4436 to 0.4644, a relative
spread of 4.7 percent, against a 20.5 percent improvement over the base model.
The difference between the best and worst architecture is therefore about a
quarter of the difference between using compensation and not using it.

We tested this directly. A Friedman test across the 32 held-out recordings
rejects the hypothesis that the seven architectures have identical error
distributions ($\chi^2 = 23.88$, $p = 5 \times 10^{-4}$), so the models are not
strictly interchangeable. However, when all 21 pairwise comparisons are made with
Wilcoxon signed-rank tests and corrected for multiple comparisons using the Holm
procedure, no pair differs significantly at the 5 percent level. The smallest
adjusted $p$ value is 0.221, for LSTM against the selective state-space model.

The appropriate conclusion is therefore narrower than architectural equivalence.
There is detectable heterogeneity among the seven, but the data do not identify
which architecture is preferable: no individual comparison survives correction,
and the spread is small relative to the effect of compensation itself. For this
task and dataset, the choice of sequence architecture is not the factor that
determines accuracy. Parameter count in particular is uninformative, since the
largest model ranks second and the smallest ranks sixth.

This has a practical consequence. Where two architectures cannot be separated on
accuracy, secondary criteria become decisive. The selective state-space model
attains comparable error with the fewest parameters but the longest training
time, owing to its sequential scan. The LSTM is the only strictly causal model in
the comparison and is therefore the natural candidate for a streaming
implementation. The Transformer achieves the lowest mean error and trains
quickly, but it processes the whole window bidirectionally and is less suited to
real-time use.

% Compile main.tex, not this file.
%==============================================================================
\section{Experiments and Results}
\label{sec:results}

%------------------------------------------------------------------------------
\subsection{Experimental Design}
\label{sec:design}

Four experiments were run, summarised in Table~\ref{tab:experiments}. The first
two correspond to the evaluation protocols of Section~\ref{sec:protocol} and
form the main results. The remaining two are ablations that remove one component
of the method at a time while holding everything else fixed. All four use the
same window length, stride, optimiser settings and epoch budget, so differences
between them are attributable to the stated change alone.

\begin{table}[!t]
\caption{Experiments. All use a window of 100 samples with stride 20, a maximum
of 100 epochs with early stopping, and a \SI{10}{\second} calibration window
unless stated otherwise.}
\label{tab:experiments}
\centering
\begin{tabular}{llp{3.0cm}}
\toprule
ID & Configuration & Question addressed \\
\midrule
E1 & Cross-day, four folds & Performance in the deployment setting \\
E2 & Within-day, three folds per session & Behaviour when between-session variation is removed \\
E3 & E1 without gain calibration & Contribution of calibration \\
E4 & E1 with a static base model & Contribution of base model dynamics \\
\bottomrule
\end{tabular}
\end{table}

All experiments were run on a single NVIDIA V100 GPU. Training the seven
architectures took 122 minutes for the four cross-day folds, 88 minutes for the
fifteen within-day folds, and 128 minutes for the dynamics ablation.

%------------------------------------------------------------------------------
\subsection{Cross-Day Performance}
\label{sec:crossday}

Table~\ref{tab:crossday} reports cross-day results. Each column is a session
held out entirely from training and validation, and each entry is the mean RMSE
over the recordings of that session. The row labelled EDM base model gives the
performance of the calibrated base model on exactly the same samples, so every
compensated result can be read against its own baseline.

\begin{table*}[!t]
\caption{Cross-day (leave-one-day-out) RMSE per held-out session. Lower is
better. The best compensated result in each column is shown in bold. Every
model improves on the base model in every fold.}
\label{tab:crossday}
\centering
\begin{tabular}{lrrrrrr}
\toprule
Model & 1st\_Comp & 2nd\_Comp & 3rd\_Base & 4th\_Valid & Overall RMSE & Overall MAE \\
\midrule
EDM base model          & 0.4884 & 1.0793 & 0.1890 & 0.3659 & 0.5581 & 0.4461 \\
\midrule
Transformer             & \textbf{0.3923} & 0.9310 & \textbf{0.1342} & \textbf{0.2450} & \textbf{0.4436} & 0.3252 \\
Conv-BiLSTM             & 0.4435 & 0.9138 & 0.1436 & 0.2480 & 0.4495 & 0.3310 \\
PINN                    & 0.4384 & 0.9063 & 0.1458 & 0.2557 & 0.4501 & 0.3363 \\
LSTM                    & 0.4536 & \textbf{0.8915} & 0.1488 & 0.2638 & 0.4520 & 0.3314 \\
Conv-BiLSTM-Transformer & 0.4482 & 0.9258 & 0.1396 & 0.2521 & 0.4547 & 0.3371 \\
Selective SSM (Mamba)   & 0.4716 & 0.9010 & 0.1494 & 0.2750 & 0.4622 & 0.3448 \\
Synergy-SSM             & 0.4637 & 0.8972 & 0.1574 & 0.2830 & 0.4644 & 0.3441 \\
\bottomrule
\end{tabular}
\end{table*}

The compensator improves on the base model in every fold and for every
architecture. Taken over all held-out recordings, the best model reduces RMSE
from 0.5581 to 0.4436, a relative reduction of 20.5 percent, and even the
weakest of the seven reaches 0.4644, a reduction of 16.8 percent. Of the 224
model and recording combinations, 175 (78 percent) show lower error than the
base model on the same samples.

The size of the improvement varies with the held-out session. It is smallest for
2nd\_Comp (17.4 percent), which has the largest force range and the highest gain
of the four, and largest for 4th\_Valid (33.0 percent) and 3rd\_Base
(29.0 percent), whose gains sit near the middle of the range spanned by the
training sessions. The improvement on 1st\_Comp is 19.7 percent. This ordering is
consistent with the interpretation developed in Section~\ref{sec:extrapolation}:
the method performs best when the held-out session is well inside the range of
conditions represented in training.

A one-sided reading of these numbers would be unwarranted without testing them,
so we treat each held-out recording as a paired observation and compare each
model against the base model with a Wilcoxon signed-rank test over the 32 test
recordings. All seven architectures improve on the base model at
$p < 4 \times 10^{-4}$, with the strongest evidence for the PINN
($p = 5.9 \times 10^{-7}$, improving on 28 of 32 recordings) and the Transformer
and Conv-BiLSTM ($p = 1.3 \times 10^{-5}$, 26 of 32 each). The benefit of
compensation is therefore consistent across recordings rather than driven by a
few favourable cases.

%------------------------------------------------------------------------------
\subsection{Within-Day Performance}
\label{sec:withinday}

Table~\ref{tab:withinday} reports the within-day results, in which training and
testing use different recordings from the same session. Across all 36 held-out
recordings the picture is markedly weaker than in the cross-day setting. Only
two of the seven architectures achieve a lower aggregate RMSE than the base
model, and the best of them improves on it by 2.9 percent, compared with
20.5 percent across sessions.

\begin{table}[!t]
\caption{Within-day RMSE, averaged over three rotating folds per session. The
base model is placed in rank order so that models falling below it are visible
directly.}
\label{tab:withinday}
\centering
\begin{tabular}{lrr}
\toprule
Model & Overall RMSE & Overall MAE \\
\midrule
PINN                    & \textbf{0.4343} & 0.3354 \\
Synergy-SSM             & 0.4467 & 0.3474 \\
\textit{EDM base model} & \textit{0.4473} & \textit{0.3562} \\
Conv-BiLSTM             & 0.4481 & 0.3513 \\
Selective SSM (Mamba)   & 0.4498 & 0.3593 \\
Transformer             & 0.4610 & 0.3645 \\
LSTM                    & 0.4624 & 0.3719 \\
Conv-BiLSTM-Transformer & 0.4633 & 0.3673 \\
\bottomrule
\end{tabular}
\end{table}

This aggregate requires careful interpretation, because it is not representative
of the individual sessions. Averaging RMSE over recordings whose force ranges
differ by more than an order of magnitude gives disproportionate weight to the
session with the largest absolute errors. Here that session is 2nd\_Comp, which
contributes nine of the 36 held-out recordings and has errors near 1.1, roughly
twenty times those of 5th\_Comp. The PINN is the only architecture that improves
on the base model on 2nd\_Comp, which is the main reason it leads the aggregate.

Examined session by session, the result is more favourable. In every one of the
five sessions at least one architecture achieves a lower RMSE than the base
model: by 7.9 percent on 1st\_Comp, 3.6 percent on 2nd\_Comp, 13.2 percent on
3rd\_Base, 13.9 percent on 4th\_Valid and 11.5 percent on 5th\_Comp. What does
not hold is consistency. The architecture that performs best differs between
sessions, and no architecture improves on the base model in all five. Over all
model and recording combinations, 168 of 252 (67 percent) improve on the base
model, compared with 78 percent in the cross-day setting.

The contrast with Section~\ref{sec:crossday} is the main finding of this
experiment, and it runs opposite to the usual expectation that within-session
evaluation is the easier case. The explanation is the amount of training data
available. A cross-day fold trains on 24 to 28 recordings, whereas a within-day
fold trains on between three and nine recordings from a single session. Removing
between-session variation also removes most of the training set, and for this
task the loss of data outweighs the gain in homogeneity. A secondary factor is
that per-session calibration (Section~\ref{sec:calibration}) already absorbs much
of the between-session gain variation before the network sees the data, so the
cross-day setting is less disadvantaged than it would otherwise be.

%------------------------------------------------------------------------------
\subsection{Ablation: Per-Session Gain Calibration}
\label{sec:ablation-calib}

% PENDING. E3 did not complete: two concurrently submitted jobs wrote to the
% same output directory, and the first to finish relocated the directory while
% E3 was still running. The run must be repeated in isolation before this
% subsection can be written.
%
% Expected content once available:
%   - cross-day RMSE for base and all seven models with calibration disabled
%   - comparison against Table IV to isolate the contribution of calibration
%   - this supports the paper's central methodological claim, so it should be
%     reported with the same per-fold detail as Section IV-B

%------------------------------------------------------------------------------
\subsection{Ablation: Base Model Dynamics}
\label{sec:ablation-lpf}

Section~\ref{sec:basemodel} argued that the first-order stage is required
because the force to sEMG relationship is not single-valued. To test this, the
cross-day experiment was repeated with the first-order stage removed, leaving a
purely static saturating nonlinearity. All other settings were unchanged.

\begin{table}[!t]
\caption{Effect of removing the first-order stage from the base model. Values
are cross-day RMSE over all held-out recordings.}
\label{tab:lpf}
\centering
\begin{tabular}{lrr}
\toprule
 & With dynamics & Static only \\
\midrule
EDM base model          & 0.5581 & 0.5683 \\
Transformer             & \textbf{0.4436} & 0.4665 \\
Conv-BiLSTM             & 0.4495 & 0.4687 \\
Conv-BiLSTM-Transformer & 0.4547 & 0.4697 \\
PINN                    & 0.4501 & 0.4845 \\
LSTM                    & 0.4520 & 0.4913 \\
Synergy-SSM             & 0.4644 & 0.4917 \\
Selective SSM (Mamba)   & 0.4622 & 0.4929 \\
\bottomrule
\end{tabular}
\end{table}

Removing the dynamics degrades both stages of the method, as shown in
Table~\ref{tab:lpf}. The base model itself worsens from 0.5581 to 0.5683, and
the best compensated result worsens from 0.4436 to 0.4665, a relative increase
in error of 4.9 percent. The degradation is consistent rather than selective:
all seven architectures perform worse with the static base model, and the number
of model and recording combinations improving on the base model falls from 175
to 164 out of 224.

Two observations follow. First, the compensator does not recover what the static
base model loses, even though it has access to a 100-sample history and could in
principle represent the missing dynamics itself. Placing the dynamics in the
physics-based component is therefore more effective than leaving them to be
learned. Second, the ranking of architectures is largely preserved between the
two conditions, which suggests that the choice of architecture and the choice of
base model structure act independently.

%------------------------------------------------------------------------------
\subsection{Transfer Boundary}
\label{sec:extrapolation}

Section~\ref{sec:protocol} explained that one session is reserved for training
because its gain lies well below that of all others. Treating it as a held-out
test session instead is informative, because it shows what happens when the
method is asked to extrapolate rather than interpolate. We therefore repeated
the cross-day experiment with that session held out, keeping all other settings
identical.

The outcome reverses the main result. The calibrated base model achieves an RMSE
of 0.0663 on this session, and every one of the seven architectures performs
worse, ranging from 0.0772 for the Transformer to 0.0838 for the Synergy-SSM.
The best model is thus 16.4 percent worse than applying no compensation at all,
and only 14 of 70 model and recording combinations improve on the base model,
compared with 78 percent in Table~\ref{tab:crossday}.

The distinguishing feature of this session is that its gain of 0.44 lies a factor
of 1.9 below the lowest gain among the remaining sessions, so no training session
brackets it from below. The compensator has learned a correction appropriate to
the range of conditions it observed and applies it outside that range, where it
is systematically wrong. This delimits the applicability of the method: it
improves accuracy when the deployment session falls inside the range of
conditions spanned by training, and degrades accuracy when it does not. Because
the gain can be estimated from the calibration segment before any prediction is
made, this condition is checkable at deployment time, and falling back to the
uncompensated base model is a safe response when it is not met.

%------------------------------------------------------------------------------
\subsection{Comparison Between Architectures}
\label{sec:architecture}

The seven architectures span convolutional, recurrent, attention-based,
state-space and physics-constrained designs, and differ by a factor of five in
parameter count, from 30{,}770 to 155{,}297. Despite this, their cross-day
results occupy a narrow band: mean RMSE ranges from 0.4436 to 0.4644, a relative
spread of 4.7 percent, against a 20.5 percent improvement over the base model.
The difference between the best and worst architecture is therefore about a
quarter of the difference between using compensation and not using it.

We tested this directly. A Friedman test across the 32 held-out recordings
rejects the hypothesis that the seven architectures have identical error
distributions ($\chi^2 = 23.88$, $p = 5 \times 10^{-4}$), so the models are not
strictly interchangeable. However, when all 21 pairwise comparisons are made with
Wilcoxon signed-rank tests and corrected for multiple comparisons using the Holm
procedure, no pair differs significantly at the 5 percent level. The smallest
adjusted $p$ value is 0.221, for LSTM against the selective state-space model.

The appropriate conclusion is therefore narrower than architectural equivalence.
There is detectable heterogeneity among the seven, but the data do not identify
which architecture is preferable: no individual comparison survives correction,
and the spread is small relative to the effect of compensation itself. For this
task and dataset, the choice of sequence architecture is not the factor that
determines accuracy. Parameter count in particular is uninformative, since the
largest model ranks second and the smallest ranks sixth.

This has a practical consequence. Where two architectures cannot be separated on
accuracy, secondary criteria become decisive. The selective state-space model
attains comparable error with the fewest parameters but the longest training
time, owing to its sequential scan. The LSTM is the only strictly causal model in
the comparison and is therefore the natural candidate for a streaming
implementation. The Transformer achieves the lowest mean error and trains
quickly, but it processes the whole window bidirectionally and is less suited to
real-time use.
